\documentclass[11pt,preview]{standalone}
%\usepackage[a4paper]{geometry}


%\usepackage[utf8]{inputenc}
\usepackage[a4paper,left=2cm,right=2cm,top=2cm,bottom=2cm]{geometry}
\usepackage{fontspec}
\usepackage[svgnames]{xcolor}
\usepackage{graphicx}
\usepackage{lipsum}
\usepackage[most]{tcolorbox}
\usepackage{tikz}
\usepackage{varwidth}
\usepackage{capt-of}
\usepackage[export]{adjustbox}
\usepackage{csquotes}
\usepackage{textgreek}
\usepackage{hyperref}
\usepackage{enumitem}
\usepackage{amsmath,amssymb}
%\usepackage{amsthm}
\usepackage{braket}
\usepackage{bbm}
\usepackage{txfonts}
\setmainfont{Times New Roman}
\usepackage{xspace}

\usepackage[style=verbose-note,date=year,bibstyle=mybibstyle,url=false,doi=false,giveninits=true]{biblatex}
  \AtEveryBibitem{%
    \clearfield{issue}}
\bibliography{references}



%#f79256 // #fbd1a2 // #7dcfb6 // #00b2ca // #1d4e89

\definecolor{c1}{HTML}{f79256} % headings
\definecolor{c2}{HTML}{fbd1a2} % links
\definecolor{c3}{HTML}{7dcfb6} % def
\definecolor{c4}{HTML}{00b2ca} % thm
\definecolor{c5}{HTML}{1d4e89} % ex


\colorlet{headingcolor}{c5}
\colorlet{solutioncolor}{c3}
\colorlet{theoremcolor}{c4}
\colorlet{definitioncolor}{c3}
\colorlet{examplecolor}{c5}
\colorlet{linkcolor}{c4!80!black}
\colorlet{recommendedcolor}{c1}

\colorlet{tagcolor}{c3!50!white}

\hypersetup{
    colorlinks=true,
    linkcolor=linkcolor,
    filecolor=linkcolor,      
    urlcolor=linkcolor,
    pdftitle={ESQC Math 2022 lecture notes},
    pdfpagemode=FullScreen,
    }

\urlstyle{rm}


%\newtheorem{myTheorem}{myTheorem}
%\theoremstyle{myDefinition}
%\newtheorem{myDefinition}{myDefinition}

\tcbset{
    theoremstyle/.style={enhanced, colback=white, colframe=blue!20, arc=0pt, 
                       before upper={\parindent15pt\noindent},
                       fonttitle=\titlefont, description color=headingcolor,  
                       colbacktitle=white, coltitle=headingcolor,    
                       top=\tcboxedtitleheight,
                       boxed title style={arc=0pt},
                       attach boxed title to top left={yshift=-\tcboxedtitleheight/2,
                                                       xshift=4mm}%
                      },
}



\newtcbtheorem{myDefinition}{Definition}{theoremstyle,colframe=definitioncolor}{def}
\newtcbtheorem{myTheorem}{Theorem}{theoremstyle,colframe=theoremcolor}{thm}
\newtcbtheorem{myLemma}{Lemma}{theoremstyle,colframe=theoremcolor}{lem}
\newtcbtheorem{myExample}{Example}{theoremstyle,colframe=examplecolor}{ex}
\newtcbtheorem{myRemark}{Remark}{theoremstyle,colframe=examplecolor}{rem}
\newtcbtheorem[no counter]{mySolution}{Solution}{theoremstyle,colframe=solutioncolor}{sol}


\newtcbtheorem[no counter]{myDefinitionx}{Definition}{theoremstyle,colframe=definitioncolor}{def}
\newtcbtheorem[no counter]{myTheoremx}{Theorem}{theoremstyle,colframe=theoremcolor}{thm}
\newtcbtheorem[no counter]{myLemmax}{Lemma}{theoremstyle,colframe=theoremcolor}{lem}
\newtcbtheorem[no counter]{myExamplex}{Example}{theoremstyle,colframe=examplecolor}{ex}
\newtcbtheorem[no counter]{myRemarkx}{Remark}{theoremstyle,colframe=examplecolor}{rem}



\newcommand{\hsc}[1]{{\authorfont\footnotesize\MakeUppercase{#1}}}

\newcommand{\mytag}[1]{\tikz[baseline=(X.base)]\node [draw=black,fill=tagcolor,semithick,rectangle,inner sep=2pt, rounded corners=3pt] (X) {\hsc{#1}};\xspace }



\usepackage{titlesec}
  \titleformat{\chapter}[hang]
    {\titlefont\huge\bfseries\color{headingcolor}}
    {}{0pt}{\huge}
%    {\thechapter}{20pt}{\huge}

  % \titleformat{\chapter}[hang]
  %   {\normalfont\huge\bfseries}
  %   {}{0pt}{\huge}

\titleformat{\section}
%  {\gothambook\Large\color{black}}
  {\titlefont\Large\bfseries\color{headingcolor}}
  {\thesection}{1em}{}
\titleformat{\subsection}
  {\titlefont\large\bfseries\color{headingcolor}}
%  {\gothambook\large\color{black}}
  {\thesubsection}{1em}{}


\newtcolorbox{defbox}{
enhanced,
boxrule=0pt,frame hidden,
borderline west={4pt}{0pt}{green!75!black},
colback=green!10!white,
sharp corners
}

%\tcbset{enhanced,colback=cyan!5!white,colframe=cyan!75!black,fonttitle=\bfseries}
\tcbset{theoremstyle}


\newcommand{\recommended}[2]{%
\begin{tcolorbox}[title={Recommended reading},sidebyside,sidebyside align=top,lower separated=false,colframe=recommendedcolor]
{#1}
\tcblower
    \includegraphics[width=\linewidth,valign=t]{#2}
\end{tcolorbox}%
}


\newcommand{\rmi}{\mathrm{i}}
\newcommand{\rmd}{\mathrm{d}}

\newcommand{\diff}[2]{\frac{\rmd #1}{\rmd #2}}
\newcommand{\pdiff}[2]{\frac{\partial #1}{\partial #2}}

\newcommand{\cl}{\operatorname{cl}}
\newcommand{\interior}{\operatorname{int}}

\renewcommand{\Re}{\operatorname{Re}}
\renewcommand{\Im}{\operatorname{Im}}
\newcommand{\RR}{\mathbb{R}}
\newcommand{\CC}{\mathbb{C}}
\newcommand{\FF}{\mathbb{F}}
\newcommand{\axiomname}[1]{\hspace{1cm} \hfill{\bf\emph{#1}}}
\newcommand{\QQ}{\mathbb{Q}}
\newcommand{\NN}{\mathbb{N}}
\newcommand{\ZZ}{\mathbb{Z}}
\newcommand{\TT}{\mathbb{T}}
\newcommand{\bvec}[1]{\mathbf{#1}}
\newcommand{\spn}{\operatorname{sp}}
\renewcommand{\epsilon}{\varepsilon}
\newcommand{\Tr}{\operatorname{Tr}}
\newcommand{\Arg}{\operatorname{Arg}}
\newcommand{\salgebra}[1]{\boldsymbol{#1}}

\newcommand{\smallO}{
  \mathchoice
    {{\scriptstyle\mathcal{O}}}% \displaystyle
    {{\scriptstyle\mathcal{O}}}% \textstyle
    {{\scriptscriptstyle\mathcal{O}}}% \scriptstyle
    {\scalebox{.7}{$\scriptscriptstyle\mathcal{O}$}}%\scriptscriptstyle
  }


\setmonofont[Scale=0.8]{DejaVu Sans Mono}
\newfontfamily{\titlefont}{Avenir Medium.otf}
\newfontfamily{\authorfont}{Avenir Light.otf}
%\DeclareTextFontCommand{\texttitlefont}{\titlefont}

\newcommand{\todo}[1]{\textcolor{red!70!black}{{\bf TODO: }#1}}



\newcommand*{\mytitle}[2]{\begingroup
\vspace*{1cm}
\begin{flushleft}\thispagestyle{empty}
{\titlefont\huge
\addfontfeature{LetterSpace=4.0}
MATHEMATICS FOR\\
QUANTUM CHEMISTRY\\
\vspace{1em}
{#1}\\
\vspace{0.5em}
{\large #2}
}
\vfill
\Large\authorfont\addfontfeature{LetterSpace=4.0}
SIMEN KVAAL \\
HYLLERAAS CENTRE\\
UNIVERSITY OF OSLO\\
simen.kvaal@kjemi.uio.no
\vfill
European Summerschool in Quantum Chemistry 2024
\end{flushleft}
\endgroup}




\begin{document}
\begin{comment}
\begin{myTheoremx}{Fundamental theorem of algebra}{}
    Every polynomial $p$ of degree $n$ over $\mathbb{C}$ have exactly $n$ roots in $\mathbb{C}$, i.e., there is a nonzero $C \in \mathbb{C}$ and $n$ numbers $r_i \in \mathbb{C}$, such that
    \[ p(z) = C (z-r_1)(z-r_2)\cdots (z-r_n). \]
\end{myTheoremx}
\end{comment}

\begin{comment}
\begin{myDefinitionx}{Euclidean space}{} Let $\FF$ be either $\RR$ or $\CC$. Let $ \FF^n$ be the set of $n$-tuples of $\FF$-numbers $\mathbf{x} = (x_1,\cdots,x_n)$, on which we define the following operations: For $\mathbf{x},\mathbf{y}\in\FF^n$ define
\begin{equation}
    \mathbf{x} + \mathbf{y} \in \FF^n, \quad (\mathbf{x}+\mathbf{y})_i = x_i + y_i \text{\axiomname{addition}},
\end{equation}
for all $1 \leq i \leq n$.
and for any $\alpha \in \FF$,
\begin{equation}
\alpha \mathbf{x} \in \FF^n, \quad (\alpha \mathbf{x})_i = \alpha x_i \text{\axiomname{scalar multiplication}}.
\end{equation}
We also define the Euclidean \emph{inner prooduct}
\begin{equation}
    \braket{\mathbf{x}, \mathbf{y}} = \bar{\mathbf{x}} \cdot \mathbf{y} = \sum_i \bar{x}_i y_i \in \FF \text{\axiomname{Euclidean inner product}}
\end{equation}
and the Euclidean norm
\begin{equation}
    \|\mathbf{y}\| = \sqrt{\braket{\mathbf{x},
    \mathbf{x}}} \in\RR.\text{\axiomname{Euclidean norm}}
\end{equation}
\end{myDefinitionx}
\end{comment}

\begin{comment}
\begin{myDefinitionx}{Standard basis}{}
The \emph{standard basis} for $\FF^n$ is the set of vectors $\{\bvec{e}_i \mid 1 \leq i \leq n\}$ such that
\begin{equation}
    (\bvec{e}_i)_j = \delta_{ij}, \axiomname{Kronecker delta symbol}
\end{equation}
i.e., 
\begin{equation}
    \bvec{e}_1 = \begin{bmatrix} 1 \\ 0 \\ \vdots \\ 0\end{bmatrix}, \quad
    \bvec{e}_2 = \begin{bmatrix} 0 \\ 1 \\ \vdots \\ 0\end{bmatrix}, \quad\text{etc.}
\end{equation}
\end{myDefinitionx}
It now follows that for every $\bvec{x}\in\FF^n$,
\begin{equation}
    \bvec{x} = \sum_{i=1}^n x_i \bvec{e}_i.
\end{equation}
It is easy to see, that 
\begin{equation}
    x_i = \braket{\bvec{e}_i,\mathbf{x}}.
\end{equation}
\end{comment}

\begin{comment}
\begin{myDefinitionx}{Linear operator}{}
    Let $A : \FF^n \to \FF^m$ be a function. We say that $A$ is \emph{linear} if it conserves the vector addition and scalar multiplication laws, i.e.,
    \begin{equation}
        A(\mathbf{x} + \mathbf{y}) = A(\mathbf{x}) + A(\mathbf{y}),
    \end{equation}
    and
    \begin{equation}
        A(\alpha\mathbf{x}) = \alpha A(\mathbf{x}).
    \end{equation}
\end{myDefinitionx}
\end{comment}

\begin{comment}
\begin{myDefinitionx}{Matrix product}{}
    Let $A \in M(n,m,\FF)$ and $B \in M(m,o,\FF)$. Then the \emph{matrix product} $C = AB \in M(n,o;\FF)$ is defined by the formula
    \begin{equation}
        C_{ik} = \sum_{j=1}^n A_{ij} B_{jk}.
    \end{equation}
    The matrix product satisfies:
    \begin{enumerate}
        \item $A(BC) = (AB)C$ \axiomname{associativity}
        \item $(A+B)C = AC + BC$ and $A(B+C) = AB + AC$ \axiomname{distributivity}
    \end{enumerate}
    However, the matrix product is \emph{not commutative}, i.e., $AB\neq BA$ in general!
 \end{myDefinitionx}
\end{comment}

\begin{comment}
\begin{myDefinitionx}{Vector space}{vector-space}
    A \emph{vector space over the field $\FF$} is a set $V$ together with a binary \emph{vector addition} $+ : V \times V \to V$ and \emph{scalar multiplication} $\cdot : \FF \times V \to V$ such that, for all $x$, $y$, $z \in V$ and all $\alpha,\beta \in \FF$, the following axioms are true:
    \begin{enumerate}
        \item There exists a $0\in V$ such that $0 + x = x$ for all $x\in V$ \axiomname{identity element for addition}
        \item  $ x + (y + z) = (x + y) + z$ \axiomname{associativity for addition}
        \item $x + y = y + x$ \axiomname{commutativity for addition}
        \item There exists $x'$ such that $x+x'=0$ \axiomname{inverse element for addition}
        \item $(\alpha\beta)\cdot x = \alpha\cdot(\beta\cdot x)$ \axiomname{compatibility of scalar and field multiplications}
        \item $1\cdot x = x$ \axiomname{identity for scalar multiplication}
        \item $(\alpha + \beta)\cdot x = \alpha\cdot x + \beta\cdot x$ \axiomname{distributivity of scalar multiplication}
        \item $\alpha \cdot (x + y) = \alpha\cdot x + \alpha \cdot y$ \axiomname{distributivity of scalar multiplication}
    \end{enumerate}
\end{myDefinitionx}
\end{comment}

\begin{comment}
\begin{myDefinitionx}{Linear independence, dimension}{}
    Let $V$ be a vector space, and $L \subset V$ a subset. The set $L$ is \emph{linearly indepdenent} if for any finite set $\{v_i \mid 1 \leq i \leq k \} \subset L$, we have
    \[ \sum_{i=1}^k a_i v_i = 0 \implies  a_i = 0 \text{ for all $i$} \]
    The \emph{dimension} of $V$ is the cardinality of the largest linearly independent subset of $V$.
\end{myDefinitionx}
\end{comment}


\begin{comment}
\begin{myDefinitionx}{Basis}{}
    Let $V$ be a vector space of finite dimension $n$. A basis is a linearly independent set of vectors $\{b_1,\cdots,b_n\}$, with exactly $n$ elements.
\end{myDefinitionx}
 
\begin{myTheoremx}{}{}
    If $B= \{b_1,\cdots,b_n\}$ is a basis for a the vector space $V$, $\dim(V)<+\infty$, then any $v \in V$ can be uniquely decomposed as
    \begin{equation}
        v = \sum_{i=1}^n v_i b_i.
    \end{equation}
\end{myTheoremx}
\end{comment}

\begin{comment}
\begin{myDefinitionx}{Inner product}{}
    Let $V$ be a vector space.
    An \emph{inner product} $\braket{\cdot, \cdot} : V\times V \to \FF$ is a map which satisfies the following axioms:
   \begin{enumerate}
       \item $\braket{x,x} \geq 0$, \quad $\braket{x,x} = 0$ if and only if $x = 0$ \hfill \axiomname{non-negative}
       \item $\braket{x,\alpha y + \beta z} = \alpha\braket{x,y} + \beta\braket{x,z}$ \axiomname{linearity}
       \item $\braket{\alpha y + \beta z, x} = \bar{\alpha}\braket{y,x} + \bar{\beta}\braket{z,x}$ \axiomname{conjugate linearity}
       \item $\braket{x,y} = \overline{\braket{y,x}}$ \axiomname{hermiticity}
   \end{enumerate}
\end{myDefinitionx}
\end{comment}

\begin{comment}
\begin{myRemarkx}{}{}
In order to study (the vector space structure of) finite dimensional Hilbert spaces, including the linear operators over these spaces, it suffices to study $\FF^n$ and matrices $M(n,m,\FF)$.
\end{myRemarkx}
\end{comment}

\begin{comment}
\begin{myDefinitionx}[{Linear subspace}{}
Let $V$ be a vector space over $\FF$. A subset $W \subset V$ is a \emph{linear subspace} if it is closed under vector addition and scalar multiplication, i.e., if
\begin{equation}
    \forall w \in W, \quad \alpha w \in W, \quad w_1 + w_2 \in W,
\end{equation}
and if $0 \in W$.
\end{myDefinitionx}
\end{comment}


\begin{comment}
\begin{myDefinitionx}[width=0.75\textwidth]{Column space}{}
    For a matrix $A = [\bvec{a}_1,\cdots,\bvec{a}_n]\in\FF^{n\times m}$ , the \emph{column space} is the set of all linear combinations of the columns $\bvec{a}_i$. This is also denoted \emph{the range} or \emph{image} of $A$, since it is the set of all vectors $A\bvec{x}$. 

    The column space is a linear vector space, written
    \begin{equation}
        \operatorname{span}\{ \bvec{a}_1, \, \bvec{a}_2, \, \cdots\, , \bvec{a}_n \}.
    \end{equation}

    The \emph{rank} of the matrix is the dimension of the column space. 

    The row space is defined similarly.

    (It is a fact that the dimension of the row space is the same as the dimension of the column space.)
\end{myDefinitionx}
\end{comment}

\begin{comment}
\begin{myDefinitionx}[width=0.75\textwidth]{Linear operators in finite dimensional spaces}{}
    Let $V$ and $W$ be  vector spaces over $\FF$ of finite dimensions $n$ and $m$, respectively. Let $\hat{A} : V \to W$ be a linear operator, i.e., $\hat{A}(v + v') = \hat{A}v + \hat{A}v'$, and $\hat{A}(\alpha v) = \alpha \hat{A}v$. We denote by $L(V,W)$ the set of such linear functions, which is a vector space.

    Let $\{v_i\} \subset V$ and $\{w_i\}\subset W$ be \emph{orthonormal bases}, i.e.,
    \begin{equation*}
        \braket{v_i,v_j}= \delta_{ij}, \quad \braket{w_i,w_j} = \delta_{ij}.
    \end{equation*}

    Then the \emph{matrix} of $\hat{A} \in L(V,W)$ relative to the given bases is
    \begin{equation*}
        A_{ij} = \braket{w_i, \hat{A} v_j}
    \end{equation*}

\end{myDefinitionx}
\end{comment}

\begin{comment}
\begin{myDefinitionx}[width=0.75\textwidth]{Unitary operator/matrix}{}
\begin{enumerate}
    %\item 
    A matrix $U \in M(n,n;\FF)$ is \emph{unitary} if, for all $\mathbf{x}$, $\mathbf{y}\in \FF^n$,
    \begin{equation}
        \braket{U\mathbf{x},U\mathbf{y}} = \braket{\bvec{x},\bvec{y}}, \quad\text{equivalently}\quad U^H U = UU^H =  \mathbbm{1}.
    \end{equation}
% \item 
%     Let $V$ be a finite-dimensional Hilbert space.
%     Let $\hat{U}\in L(V)$ be a linear operator.

%     $\hat{U}$ is \emph{unitary} if, for all $u,v\in V$,
%     \begin{equation*}
%         \braket{\hat{U} u, \hat{U} v} = \braket{u, v}, \quad \text{equivalently} \quad \hat{U}^\dag \hat{U} = \hat{U}\hat{U}^\dag = \mathbbm{1}.
%     \end{equation*}
\end{enumerate}
\end{myDefinitionx}
\end{comment}

\begin{comment}
\begin{myDefinitionx}[width=0.75\textwidth]{Hermitian operator/matrix}{}
%\begin{enumerate}
    %\item 
    A matrix $A \in M(n,n;\FF)$ is \emph{Hermitian} if, for all $\mathbf{x}$, $\mathbf{y}\in \FF^n$,
    \begin{equation}
        \braket{\mathbf{x},A\mathbf{y}} = \braket{A \bvec{x},\bvec{y}}, \quad\text{equivalently}\quad A^H = A.
    \end{equation}
% \item 
%     Let $V$ be a finite-dimensional Hilbert space.
%     Let $\hat{A}\in L(V)$ be a linear operator.

%     $\hat{A}$ is \emph{Hermitian} if, for all $u,v\in V$,
%     \begin{equation*}
%         \braket{ u, \hat{A} v} = \braket{\hat{A} u, v}, \quad \text{equivalently} \quad \hat{A}^\dag = \hat{A}.
%     \end{equation*}
%\end{enumerate}
\end{myDefinitionx}
\end{comment}

\begin{comment}
\begin{myTheoremx}[width=0.65\textwidth]{Spectral theorem for Hermitian operators}{}
    Suppose $A \in \FF^{n\times n}$ is Hermitian, i.e., $A^H = A$. Then, there exists an orthonormal basis $\{\mathbf{u}_1,\cdots,\mathbf{u}_n\}$, and real numbers $\{\lambda_1,\cdots,\lambda_n\}$, such that
    \begin{equation*}
        H \mathbf{u}_i = \lambda_i \mathbf{u}_i
    \end{equation*}   
    Equivalently, 
    \begin{equation*}
        A = \sum_{i=1}^n \mathbf{u}_i \lambda_i \mathbf{u}_i^H = U \Lambda U^H
    \end{equation*}
    where $\mathbf{u}_i$ is the $i$th column of $U$, and where $\Lambda$ is a diagonal matrix with elements $\Lambda_{ij} = \lambda_i \delta_{ij}$.
\end{myTheoremx}
\end{comment}

\begin{comment}
\begin{myTheoremx}[width=0.65\textwidth]{Singular value decomposition}{}
    Let $A \in M(n,m,\FF)$ be a matrix, and let $k = \min(n,m)$. There exists $k$ \emph{singular values} $\sigma_i \geq 0$ and $k$ \emph{left singular vectors} $\bvec{u}_i$, and $k$ \emph{right singular vectors} $\bvec{v}_i$, such that
    \begin{equation*}
        A = \sum_{i=1}^k \bvec{u}_i \sigma_i \bvec{v}_i^H = U \Sigma V^H,
    \end{equation*}
    where $U = [\bvec{u}_1,\cdots,\bvec{u}_k]$, $V = [\bvec{v}_1,\cdots,\bvec{v}_k]$, $\Sigma = \operatorname{diag}(\sigma_1,\cdots,\sigma_k)$. Equivalently,
    \begin{equation*}
        A \mathbf{v}_i = \sigma_i \mathbf{u}_i.
    \end{equation*}
    The rank of $A$ is the number of nonzero singular values. The decomposition is unique if all the singular values are distinct.
\end{myTheoremx}
\end{comment}





\begin{comment}
\begin{myDefinitionx}[width=0.65\textwidth]{Topologically important sets}{}
    \begin{enumerate}
%        \item 
%Let $\mathbf{x} \in \RR^n$ and $\epsilon>0$ be given. The set $B_\epsilon(x) = \{ y \in M \mid d(x,y) < \epsilon \}$ is called an \emph{$\epsilon$-ball around $x$}.
    \item A subset $S \subset \RR^n$ is called \emph{open} if, for every $\bvec{x} \in S$, there is an $\epsilon > 0$ such that $B_\epsilon(\bvec{x}) \subset S$.
%    \item A \emph{neighborhood} of $x \in S$ is a subset $N$ such that some $\epsilon$-ball around $x$ is a subset of $N$. 
    \item A subset $S$ is called \emph{closed}  if $S^\complement =  \RR^n \setminus S$ is open.
    \item The \emph{closure} $\cl(S)$ is the smallest closed set that contains $S$.
    \item The \emph{interior} $\interior(S)$ is the set of all those $\bvec{x}\in S$ around which there exists an $\epsilon$-ball in $S$
    \item The \emph{boundary} $\partial S$ is the intersection $\cl(S^\complement) \cap \cl(S) = S \setminus \interior(S)$
%    \item A subset $S \subset \RR^n$ is \emph{dense} if its closure is all of $\RR^n$
    \end{enumerate}
\end{myDefinitionx}
\end{comment}

\begin{comment}
\begin{myDefinitionx}[width=0.65\textwidth]{Limit}{}
    Let $f : \Omega \subset \RR^n \to \RR^m$, where $\Omega$ is open. Let $\bvec{x}_0\in \Omega \cup \partial\Omega$, and let $N$ be a neighborhood of $\bvec{b}\in \RR^m$.

    We say that $f$ is \emph{eventually in $N$ as $\bvec{x}$ approaches $\bvec{x}_0$}, if there exists a neighborhood $U$ of $\bvec{x}_0$, such that $\bvec{x}
    \in U$ but $\bvec{x}\neq\bvec{x}_0$ and $\bvec{x}\in \Omega$ imply $f(x)\in N$.

    We say that \emph{$f(\bvec{x})$ approaches $\bvec{b}$ as $\bvec{x}$ approaches $\bvec{x}_0$},
    \begin{equation}
        \lim_{\bvec{x}\to\bvec{x}_0} f(\bvec{x}) = \bvec{b} \quad\text{or}\quad f(\bvec{x}) \to \bvec{b} \; \text{as} \; \bvec{x}\to\bvec{x}_0,
    \end{equation}
    when, given any neighborhood $N$ of $\bvec{b}$, $f$ is eventually in $N$ as $\bvec{x}$ approaches $\bvec{x}_0$.
\end{myDefinitionx}
\end{comment}

\begin{comment}
\begin{myDefinitionx}[width=0.5\textwidth]{Continuity}{}
    Let $f : \Omega \subset \RR^n\to\RR^m$. Let $\bvec{x}_0\in\Omega$. We say that $f$ is \emph{continuous at $\bvec{x}_0$} if
    \begin{equation*}
        \lim_{\bvec{x}\to\bvec{x}_0} f(\bvec{x}) = f(\bvec{x}_0).
    \end{equation*}
\end{myDefinitionx}
\end{comment}

\begin{comment}
\begin{myTheoremx}[width=0.65\textwidth]{Properties of continuous functions}{}
Let $f, g : \Omega\subset\RR^n \to \RR^m$ be functions with a common domain $\Omega$, continuous at $\bvec{x}_0$:
Then:
\begin{enumerate}
    \item $f+g$ and $\alpha f$ for any $\alpha\in\RR$ are continuous at $\bvec{x}_0$.
    \item In the scalar-valued case $m=1$, the product $fg$ is continuous at $\bvec{x}_0$
    \item If $f \neq 0$ in all of $\Omega$, then $1/f$ is continuous at $\bvec{x}_0$
    \item The component functions $f_i : \Omega \to \RR$ are all continuous at $\bvec{x}_0$. The converse is also true.
\end{enumerate}  
\end{myTheoremx}
\end{comment}

\begin{comment}
\begin{myTheoremx}[width=0.65\textwidth]{Compositions of functions}{}
Let $f : \Omega \subset \RR^n \to \RR^m$ be continuous at $\bvec{x}_0\in \Omega$, and $g : \Omega' \subset \RR^m \to \RR^o$. Suppose $f[\Omega] \subset \Omega'$, and let $g$ be continuous at $\bvec{y}_0 = f(\bvec{x}_0)$. Then $h : \Omega\subset \RR^n\to\RR^o$,
\[ h(\bvec{x}) = g(f(\bvec{x}_0) \]
is continuous at $\bvec{x}_0$.
\end{myTheoremx}
\end{comment}


\begin{comment}
\begin{myDefinitionx}[width=0.65\textwidth]{Partial derivative}{}
Let $f : \Omega\subset\RR^n \to \RR$ be a scalar-valued function, $\Omega$ open. The \emph{partial derivatives} with respect to the variable $x_i$ are defined by
\begin{equation*}
    \pdiff{}{x_i} f(\bvec{x}) = \lim_{h\to 0} \frac{f(\bvec{x} + h\bvec{e}_i) - f(\vec{x}) }{h}
\end{equation*}
if the limit exists. 

In the case $f : \Omega\subset\RR^n \to \RR^m$, the the partial derivatives are defined componentwise, i.e.,
\begin{equation*}
    \pdiff{}{x_i} f_j(\bvec{x}).
\end{equation*}
\end{myDefinitionx}
\end{comment}

\begin{comment}
\begin{myExamplex}[width=0.5\textwidth]{}{}
    let $f : \RR^2 \to \RR$, $(x,y)\mapsto x^{1/3} y^{1/3}$. Since $f(x,0) = 0$ and $f(0,y) = 0$, 
    \begin{equation}
        \pdiff{}{x} f(0,0) = \pdiff{}{y}f(0,0) = 0.
    \end{equation}
    But along the ``diagonal''
    \begin{equation}
        g(x) = f(x,x) = x^{2/3}.
    \end{equation}
    The derivative of $g(x)$ is
    \begin{equation}
        g'(x) = \frac{2}{3} x^{-1/3} \to +\infty \quad \text{as $x\to 0$}
    \end{equation}
\end{myExamplex}
\end{comment}

\begin{comment}
\begin{minipage}{0.65\textwidth}
\begin{myDefinitionx}{Differentiable}{}
    Let $f : \Omega \subset \RR^n \to \RR^m$, with $\Omega$ open. We say that $f$ is \emph{differentiable at $\bvec{x}_0\in\Omega$} if the partial derivatives all exist at $\bvec{x}_0$, and if
    \begin{equation*}
        \lim_{\bvec{x}\to\bvec{x}_0} \frac{\| f(\bvec{x}) - f(\bvec{x}_0) - M(\bvec{x}-\bvec{x}_0)\|}{\|\bvec{x}-\bvec{x}_0\|} = 0,
    \end{equation*}
    where $M = Df(\bvec{x}_0)$, the \emph{derivative}, is the matrix of partial derivatives, 
    \begin{equation*}
        M_{ij} = \frac{\partial f_i(\bvec{x}_0)}{\partial x_j}.
    \end{equation*}
    and where $M(\bvec{x}-\bvec{x}_0)$ is the matrix-vector product applied to $\bvec{x}-\bvec{x}_0$.
\end{myDefinitionx}
\end{minipage}
\end{comment}

\begin{comment}
\begin{myTheorem}[width=0.65\textwidth]{}{}
    If $f$ is differetiable at $\bvec{x}_0$, it is continuous at $\bvec{x}_0$.
\end{myTheorem}


\begin{myTheorem}[width=0.65\textwidth]{Condition for differentiability}{}
    Let $f : \Omega \subset \RR^n \to \RR^m$, with $\Omega$ open. Suppose the partial derivatives all exist at $\bvec{x}_0$, and furthermore that they are all continuous \emph{in a neighborhood} of $\bvec{x}_0$. Then $f$ is differentiable at $\bvec{x}_0$.
\end{myTheorem}
\end{comment}

\begin{comment}
\begin{myDefinition}[width=0.5\textwidth]{$C^1$ functions}{}
        A function whose partial derivatives exist and are continuous throughout its open domain $\Omega$ is said to be of class $C^1$.
\end{myDefinition}
\end{comment}

\begin{comment}
\begin{myTheoremx}{Properties of the derivative}{}
\begin{enumerate}
    \item Let $f : \Omega \subset \RR^n \to \RR^m$ be differentiable at $\bvec{x}_0\in\Omega$, and let $c \in \RR$. Then $h(\bvec{x}) = c f(\bvec{x})$ is differentiable at $\bvec{x}_0$, and
    \begin{equation}
        D h(\bvec{x}_0) = c D f(\bvec{x}_0).
    \end{equation}
    \item Let $g : \Omega \subset \RR^n \to \RR^m$ be another function differentiable at $\bvec{x}_0$. Then $h(\bvec{x}) = f(\bvec{x}) + g(\bvec{x})$ is differentiable at $\bvec{x}_0$, and
    \begin{equation}
        D h(\bvec{x}_0) = D f(\bvec{x}_0) + Dg(\bvec{x}_0).
    \end{equation}
    \item\label{prodrule} Let $f,g : \Omega \subset\RR^n\to \RR$ be \emph{scalar-valued} functions, differentiable at $\bvec{x}_0\in\Omega$. Then $h(\bvec{x}) = f(\bvec{x})g(\bvec{x})$ is differentiable at $\bvec{x})_0$, and
    \begin{equation}
        Dh(\bvec{x}_0) = g(\bvec{x}_0)Df(\bvec{x}_0) + f(\bvec{x}_0) Dg(\bvec{x}_0).
    \end{equation}
    \item As in \ref{prodrule}, and additionally that $g >0$ thrughout $\Omega$. Then $h(\bvec{x}_0) = f(\bvec{x}_0)/g(\bvec{x}_0)$ is differentiable at $\bvec{x}_0$, and
    \begin{equation}
        Dh(\bvec{x}_0) = \frac{g(\bvec{x}_0)Df(\bvec{x}_0) - f(\bvec{x}_0)} {[g(\bvec{x}_0)]^2}
    \end{equation}
\end{enumerate}
\end{myTheoremx}
\end{comment}

\begin{comment}
    
\begin{myTheoremx}[width=0.62\textwidth]{Chain rule}{}
    Let $\Omega \subset \RR^n$ and $\Omega' \subset \RR^m$ be open sets, and let $g : \Omega \to \RR^m$ with $g[\Omega] \subset \Omega'$. Let $f : \Omega' \to \RR^o$. Thus, $h = f \circ g : \Omega \to\RR^o$ is defined. Suppose $g$ is differentiable at $\bvec{x}_0\in\Omega$, and $f$ is differentiable at $\bvec{y}_0 = f(\bvec{x}_0)\in\Omega'$. Then  $f \circ h$ is differentiable at $\bvec{x}_0$ with derivative
    \begin{equation*}
        D(f \circ g)(\bvec{x}_0) = Df(\bvec{y}_0)Df(\bvec{x}_0),
    \end{equation*}
    i.e., the matrix product of the Jacobian matrices.
\end{myTheoremx}

\end{comment}

\begin{comment}
\begin{myExamplex}[width=0.62\textwidth]{Function along a path}{}
Let $\bvec{c} : \RR \to \RR^3$ be a path tracing out a curve in space, say the flight of a drone with a temperature sensor. The position $\bvec{c}(t) = [x(t),y(t),z(t)]^T$ is the position of the drone at time $f$. Let $f : \RR^3 \to \RR$ be a function, say temperature, in space. The temperature registered by the drone as function of time is $h(t) = (g \circ \mathbf{c})(t) = g(\mathbf{c}(t))$. The rate of change of $h(t)$ is
\begin{equation*}
    \diff{h}{t} = \pdiff{h}{x}\pdiff{x}{t} + \pdiff{h}{y}\pdiff{y}{t} + \pdiff{h}{z}\pdiff{z}{t}.
\end{equation*}
\end{myExamplex}
\end{comment}

\begin{comment}
\begin{myTheoremx}[width=0.62\textwidth]{Second-order Taylor formula}
    Let $f : \Omega \subset \RR^n \to \RR$ be of class $C^2$. Then we may write
    \begin{equation*}
        f(\bvec{x}_0 + \bvec{h}) = f(\bvec{x}_0) + Df(\bvec{x}_0)\bvec{h} + \frac{1}{2} \bvec{h}^T D^2 f(\bvec{x}_0) \bvec{h} + R_2(\bvec{h},\bvec{x}_0),
    \end{equation*}
    where the \emph{remainder} satisfies $R_2(\bvec{h},\bvec{x}_0)/\|\bvec{h}\|^2 \to 0$ as $\bvec{h}\to 0$, written
    \begin{equation*}
    R_2(\bvec{h},\bvec{x}_0) = \smallO(\|\bvec{h}\|^2).
    \end{equation*}
    The symbol $D^2 f(\bvec{x}_0)$ is the \emph{Hessian} of $f$, the matrix of second-order mixed partial derivatives, a symmetric matrix.
\end{myTheoremx}
\end{comment}

\begin{comment}
\begin{myExamplex}[width=.65\textwidth]{}{}
    Compute the second-order Taylor polynomial of \\$f(x,y) = \exp(-x^2 - y^2)$ at $(0,0)$.

    \begin{equation}
    Df(x,y) = [-2x f(x,y), -2y f(x,y)],
    \end{equation}
    \begin{equation} D^2 f(x,y) = \begin{bmatrix} (4x^2 -2)f(x,y) & 4xy f(x,y) \\ 4xy f(x,y) & (4y^2 - 2)f(x,y) \end{bmatrix}
    \end{equation}
    \begin{equation}
    f(0,0) = 1, \quad Df(0,0) = [0,0], \quad D^2 f(0,0) = \begin{bmatrix} -2 & 0 \\ 0 & -2 \end{bmatrix}
    \end{equation}
    \begin{equation}
        f(x,y) = 1 - (x^2 + y^2) + \smallO(x^2 + y^2).
    \end{equation}
\end{myExamplex}
\end{comment}

\begin{comment}
\begin{myTheoremx}[width=0.7\textwidth]{Classification of critical points} 
Let
$f : \Omega \subset \RR^n \to \RR$, with $\Omega$ being an open domain. Let $f$ be of class $C^2$. Let $H = D^2f(\bvec{x})$ be the second derivative (Hessian) at a critical point $\bvec{x}\in\Omega$, i.e., $Df(\bvec{x})=0$. 
Then we have:
\begin{enumerate}
    \item
     If all the eigenvalues of $H$ are positive, then $\bvec{x}$ is a local minimium.
     \item
     If all the eigenvalues of $H$ are negative, then $\bvec{x}$ is a local maximum.
     \item If there are eigenvalues of $H$ with both positive and negative values, but no zero eigenvalues, then $\bvec{x}$ is a saddle point.
     \item If some eigenvalues are zero, we cannot conclude based on second-order Taylor polynomials.
\end{enumerate}
\end{myTheoremx}
\end{comment}

\begin{comment}
\begin{myDefinitionx}[width=0.65\textwidth]{Complex number operations}{}
    Let $z = x + \rmi y \in \CC$.
    \begin{itemize}
        \item $\Re z = x$,  $\Im z = y$ \axiomname{real and imaginary part}        \item $\bar{z} = z^* = x - \rmi y$ \axiomname{complex conjugate}
        \item $z = r e^{i\theta}$, \axiomname{polar form}\\
        where
$ e^{\rmi\theta} = \cos\theta + \rmi\sin\theta $ \axiomname{Euler's formula}
        \item $\Arg z = \theta$ \axiomname{argument/angle/phase}
        \item $|z|^2 = \bar{z}z = \Re z^2 + \Im z^2 = r^2$ \axiomname{squared modulus/norm}
    \end{itemize}
\end{myDefinitionx}
\end{comment}

\begin{comment}
\begin{myDefinitionx}[width=0.85\textwidth]{Complex differentiability}{}
    The function $f : U \to \CC$, $U \subset^{\text{open}} \CC$, is (complex) differentiable at $z \in U$ if the limit
    \begin{equation}
        \lim_{h \to 0} \frac{f(z+h) - f(z)}{h} = f'(z) = \diff{f}{z}
    \end{equation}
    exists. The expression $h\to 0$ means the same as in the $\RR^2$ case.

    If $D$ is an open domain in $\CC$, and if $f(z)$ is complex differentiable for all $z \in D$, we say that $f$ is \emph{analytic in $D$}.
\end{myDefinitionx}
\end{comment}

\begin{comment}
\begin{myExamplex}[width=0.7\textwidth]{Derivative of monomial}{}
Let us apply the definition of the derivative to $f(z) = z^n$.
\begin{equation}
    f(z+h) = (z+h)^n = z^n + h n z^{n-1} + \text{higher order terms}.
\end{equation}
Thus
\begin{equation}
    \frac{f(z+h)-f(z)}{h} = \frac{h n z^{n-1} + \text{h.o.t.}}{h} = n z^{n-1} + \text{h.o.t.},
\end{equation}
so that the limit becomes
\begin{equation}
    \diff{}{z} z^{n} n z^{n-1}.
\end{equation}
We were able to perform the limit just by doing complex algebra. Notably, $z\in \CC$ was completely arbitrary, so the derivative exists everywhere.
\end{myExamplex}
\end{comment}

\begin{comment}
\begin{myExamplex}[width=0.7\textwidth]{Derivative of $\bar{z}$ does not exist}{}
Let us try to see if $f(z) = \bar{z}$ is differentiable. Let us consider the limit $h = \delta x \to 0$ in $\RR$. 
\begin{equation}
    \lim_{\delta x\to 0} \frac{f(x + \delta x + \rmi y) - f(x + \rmi y)}{\delta x} = \frac{\delta x}{\delta x} = 1.
\end{equation}
However, if we allow $h = \rmi \delta y \to 0$ instad, with $\delta y\in \RR$, then
\begin{equation}
    \lim_{\delta y\to 0} \frac{f(x + \rmi (\delta y + y)) - f(x + \rmi y)}{\delta y} = \frac{-\rmi \delta y}{\delta y} = -\rmi.
\end{equation}
Since the two limits are not the same, the complex limit cannot exist, since limits are unique.

This is in fact a very simple example of a continuous function from $\CC\to\CC$ which is not differentiable anywhere! Such an exampe is much harder to find for functions $\RR^2\to\RR^2$.
\end{myExamplex}
\end{comment}

\begin{comment}
\begin{myDefinitionx}[width=0.6\textwidth]{Series}{}
Given a sequence $(z_n)$ of complex numbers, we define the \emph{partial sums}
\begin{equation}
    S_N = \sum_{n=0}^N z_n.
\end{equation}
If the partual sums converge, $\lim S_N = S \in \CC$, then we denote that sum by
\begin{equation}
    S = \sum_{n=0}^\infty z_n = z_0 + z_1 + z_2 + \cdots
\end{equation}
\end{myDefinitionx}
\end{comment}

\begin{comment}
\begin{myExamplex}[width=0.7\textwidth]{Geometric series}{}
The geometric series,
\begin{equation}
   f(z) =  \frac{1}{1-z} = 1 + z + z^2 + \cdots \qquad \text{for $|z| < 1$}.
\end{equation}
The function $f(z)$ is complex differentiable at any $z \neq 1$:
\begin{equation}
    f(z + h) = \frac{1}{1-z-h} = \frac{1}{1-z} \frac{1}{1 - h/(1-z)} = \frac{1}{1-z}(1 + \frac{h}{1-z} + \text{h.o.t.}),
\end{equation}
so that $f'(z) = \frac{1}{(1-z)^2}$.

We note that $f$ is \emph{divergent} as $z\to 1$, this is an example of a \emph{pole} of $f$.
\end{myExamplex}
\end{comment}

\begin{comment}

\begin{myTheoremx}[width=0.8\textwidth]{Power series}{}
A \emph{power series} is a series of the type
\begin{equation}
    \sum_{n=0}^\infty a_n z^n.
\end{equation}
The \emph{radius of convergence} of the power series is given by
\begin{equation}
    \frac{1}{R} = \limsup_{n\to\infty} |a_n|^{1/n}.
\end{equation}
In particular, $R > 0$ if $\{|a_n|^{1/n}\}$ stays bounded as $n$ grows.

A power series is complex differentiable inside the radius of convergence, and the derivative is computed term by term,
\begin{equation}
    \diff{}{z} \sum_{n=0}^\infty a_n z^n = \sum_{n=1}^\infty n a_n z^{n-1}.
\end{equation}
The radius of convergence of the derivative is again $R$. It follows that a power series is \emph{infinitely differentiable}.
\end{myTheoremx}
\end{comment}



\begin{comment}
\begin{myDefinitionx}[width=0.6\textwidth]{Complex line integral}{}
Let $f : D \to \CC$ be continuous, and let $\Gamma$ be a (piecewise) smooth oriented curve parameterized by $\gamma : I \to \CC$.
The complex line integral of $f$ along $\Gamma$ is now defined as
\begin{equation}
    \int_\Gamma f(z) \, \rmd z = \int_I f(\gamma(t)) \gamma'(t) \, \rmd t,
\end{equation}
which is independent of parameterization. Note that $\rmd z = \gamma'(t)\rmd t$, an infinitesimally small piece of the curve. 
\end{myDefinitionx}
\end{comment}

\begin{comment}
\begin{myTheoremx}[width=0.6\textwidth]{Cauchy theorem}{}
Let $f : D \to \CC$, where $D$ is a simply connected open domain. Let $\Gamma$ be a piecewise smooth simple closed curve in $D$. Then,
\begin{equation}
    \oint_\Gamma f(z)\, \rmd z = 0.
\end{equation}
\end{myTheoremx}
\end{comment}

\begin{comment}
\begin{myTheoremx}[width=0.65\textwidth]{Cauchy integral formula}{}
    Let the function $f : D \to \CC$ be complex differentiable, and let $\Gamma$ be a simple closed curve in $D$. Then,
    \begin{equation}
    f(z_0) = \frac{1}{2\pi \rmi} \oint_\Gamma \frac{f(z)}{z - z_0} \, \rmd z.
    \end{equation}
\end{myTheoremx}
\end{comment}

\begin{myTheoremx}[width=0.64\textwidth]{}{}
     Let $D$ be simply connected, and let $f : D \to \CC$ be complex analytic in $D$. Then $f$ is \emph{infinitely} many times differentiable, and we have the power series representation
     \begin{equation}
        f(z) = \sum_{n=0}^\infty a_n(z-z_0)^n, \quad \text{where} \quad a_n = \frac{f^{(n)}(z)}{n!} 
     \end{equation}
     The derivatives are given by the formula
     \begin{equation}
         f^{(n)}(z) = \frac{1}{2\pi\rmi} \oint_\Gamma \frac{f(w)}{(w-z)^{n+1}} \, \rmd w.
     \end{equation}
\end{myTheoremx}

\end{document}